\subsection{What the quintic weight bounds give in the cyclotomic family}
\label{sec:kt-quintic-iwasawa-specialization}

This comparison uses the actual inverse limit of class groups.  It
does not construct another abstract norm system.  Its conclusion is
that the existing prime-divisible-weight bound adds no further
condition to the cyclic presentation obtained from the base class
bound.  No order-five case is removed.

Let $p\equiv1\pmod {10}$ be prime, let $F$ be the cyclic quintic
field of conductor $p$, and write $F_\infty=\bigcup_nF_n$ for its
cyclotomic $\mathbf Z_p$-extension.  Put
\[
 \Delta_5=\operatorname{Gal}(F/\mathbf Q),\quad
 \Gamma=\operatorname{Gal}(F_\infty/F),\quad
 A_n=\operatorname{Cl}(F_n)\{p\},\quad
 X=\varprojlim_{N}A_n,\quad
 \Lambda=\mathbf Z_p[[\Gamma]]=\mathbf Z_p[[T]].
 \tag{QI.1}
\]
Choose $\gamma\in\Gamma$ whose pro-$p$ cyclotomic value is $1+p$,
and put $T=\gamma-1$.  The action of $\Delta_5$ on $X$ commutes
with $\Gamma$.  Since $p\equiv1\pmod5$, all its characters have
integral Teichm\"uller values in $\mathbf Z_p$.

\begin{proposition}[KT-QUINTIC-IWASAWA-SPECIALIZATION-SCOPE]
For each nontrivial character $\alpha$ of $\Delta_5$, there is a
canonical norm-induced isomorphism
\[
                    X^\alpha/TX^\alpha\simeq A_0^\alpha.
 \tag{QI.2}
\]
Under KT-ORDER-FIVE-LENGTH-ONE, either $X^\alpha=0$, or
\[
       X^\alpha\simeq\Lambda/I_\alpha,\qquad
       \{f(0):f\in I_\alpha\}=p\mathbf Z_p.             \tag{QI.3}
\]
The second case occurs precisely when $A_0^\alpha\simeq\mathbf F_p$.

For every positive even $e$ with $\omega^{-e}=\alpha$, put
$t_e=(1+p)^{-e}-1$.  The strict group of (O5S.1) satisfies
\[
 \mathcal S_p(e)^\vee\simeq
       X^\alpha/(T-t_e)X^\alpha.
 \tag{QI.4}
\]
In the nonzero case of (QI.3), this is
\[
 \mathcal S_p(e)^\vee\simeq
   \mathbf Z_p/\{f(t_e):f\in I_\alpha\}.                \tag{QI.5}
\]
In particular, if $p\mid e$, the ideal in (QI.5) is exactly
$p\mathbf Z_p$.  Thus the already proved exponent-one conclusion
in those weights is automatic from (QI.3).
\end{proposition}

\begin{proof}
The unique prime above $p$ is totally ramified in every
$F_n/F_m$ and is principal in every layer, as proved in (O5G.4).
Let $H_m$ be the Hilbert $p$-class field of $F_m$.
Total ramification gives $H_m\cap F_n=F_m$.  Class field theory
therefore makes $N_{F_n/F_m}:A_n\to A_m$ surjective: under
Artin reciprocity it is restriction to the unramified subextension
$H_mF_n/F_n$ followed by the natural identification with
$\operatorname{Gal}(H_m/F_m)$.

The genus-field calculation (O5G.6) identifies the maximal
subextension of $H_n$ abelian over $F$ with $F_nH_0$.
Consequently the same norm induces an isomorphism
\[
                 A_n/TA_n\xrightarrow{\ \sim\ }A_0.
 \tag{QI.6}
\]
One can also obtain its injectivity by comparing the already
proved equality $|A_n^{\Gamma}|=|A_0|$ with norm surjectivity:
kernel and cokernel of $T$ on the finite group $A_n$ have equal
orders.  These isomorphisms commute with norms and with $\Delta_5$.

Taking inverse limits of (QI.6) gives $X/TX\simeq A_0$.
Here there is no closure ambiguity.  The norm maps on $TA_n$
are surjective.  Moreover $\varprojlim TA_n=TX$: for a compatible
sequence $z_n\in TA_n$, the finite nonempty sets of solutions
of $Tx_n=z_n$ have a nonempty inverse limit by compactness.
Equivalently, every finite collection of compatibility equations
can be solved by choosing a solution at its largest level.
The continuous image $TX$ of the compact group $X$ is closed.
Integral character idempotents prove (QI.2).

If $A_0^\alpha=0$, compact Nakayama gives $X^\alpha=0$.
If $A_0^\alpha=\mathbf F_p$, lift its generator to $x\in X^\alpha$.
The image $\Lambda x$ is compact and closed.  The quotient
$Y=X^\alpha/\Lambda x$ satisfies $Y/(p,T)Y=0$ and must vanish.
For clarity, any nonzero such compact quotient has a nonzero
finite continuous quotient inherited from some $A_n$; on that
finite quotient the ideal $(p,T)$ is nilpotent, since $p$ is
nilpotent and $(1+T)^{p^n}=1$.  Surjectivity of $(p,T)$ on
the quotient would be impossible.  This proves cyclicity without
an additional finite-generation assumption.  If $I_\alpha$ is
the kernel of $\Lambda\to X^\alpha$, reducing modulo $T$ in
(QI.2) gives exactly the constant-term assertion (QI.3).

We next prove the strict control statement rather than assuming
that all weights use the same finite-level class group.
Let $D_e=(\mathbf Q_p/\mathbf Z_p)(\chi_p^{-e})$ and
$\mathcal G=\operatorname{Gal}(F_\infty/\mathbf Q)
 =\Delta_5\times\Gamma$.  The coefficient $D_e$ is trivial on
$G_{F_\infty}$: the remaining finite cyclotomic subgroup has
order $(p-1)/5$ and is killed by $\chi_p^{-e}$.
The finite factor $\Delta_5$ acts by $\alpha$ on the cohomology
of $\Gamma$ with these coefficients.  A nontrivial value of
$\alpha-1$ is a $p$-adic unit.  Averaging over $\Delta_5$ thus
gives
\[
                  H^i(\mathcal G,D_e)=0\qquad(i\ge0).
 \tag{QI.7}
\]
The same assertion holds for the local cyclotomic quotient at
$p$, which is the same group $\mathcal G$ because there is only
one prime above $p$ and it is totally ramified.

Inflation--restriction now identifies global $H^1$ with the
$\mathcal G$-equivariant continuous homomorphisms from
$G_{F_\infty}$ to $D_e$, subject to ramification only above $p$.
The local restriction map is injective on the corresponding
local quotient, by the degree-one part of (QI.7).  Therefore the
strict condition is exactly that the homomorphism vanish on
the decomposition group at the unique $p$-prime.
It follows that the strict group is
\[
          \operatorname{Hom}_{\mathcal G,\mathrm{cont}}(X,D_e).
 \tag{QI.8}
\]
We use here the standard class-field inverse-limit description
of $X$ as the maximal abelian unramified pro-$p$ Galois group.
Its variant requiring complete splitting at $p$ is identical
in this tower: at every finite level the $p$-prime is principal,
so its Frobenius class in $A_n$ is zero.  One may equivalently
take the inverse limits of the $p$-split class groups, which
are the same $A_n$.  A finite unramified extension of the tower
in which the $p$-prime splits descends to such an extension at
a sufficiently high finite level.  The local splitting lets
one enlarge the level to remove the finite local extension
at $p$; unramified base change away from $p$ cannot remove
nontrivial inertia.  This is the usual finite-level descent
in that class-field description.

In (QI.8), equivariance on $\Delta_5$ selects $X^\alpha$.
Equivariance on $\Gamma$ says
$f(\gamma x)=(1+p)^{-e}f(x)$, so it selects the quotient by
$T-t_e$, with precisely the sign displayed in (QI.4).
Pontryagin duality proves (QI.4) and (QI.5).
The class-field convention is the same one as (O5S.9), where
the source class character is $\omega^{-e}$.

Finally, consider any ideal $I\subset\mathbf Z_p[[T]]$ whose
constant-term image is $p\mathbf Z_p$.  Every $f\in I$ has
$f(0)\in p\mathbf Z_p$, and at least one has
$f(0)\in p\mathbf Z_p^\times$.  For $t\in p^2\mathbf Z_p$,
\[
              f(t)\equiv f(0)\pmod {p^2}
              \quad\Longrightarrow\quad
              \{f(t):f\in I\}=p\mathbf Z_p.              \tag{QI.9}
\]
The lifting-the-exponent formula gives
$v_p(t_e)=1+v_p(e)$, so (QI.9) applies whenever $p\mid e$.
This completes the proof.
\end{proof}

\paragraph{What this leaves unevaluated.}
The cyclic presentation in (QI.3) already follows from the base
class bound and the one-prime norm calculation.  All integral
weight specializations of a cyclic $\Lambda$-module are cyclic.
Equation (QI.9) shows that the extra exponent bound on the entire
$p$-divisible weight disk supplies no additional equation on
$I_\alpha$.  The proof gives neither a bound on the Iwasawa
$\lambda$-invariant nor finiteness of $X^\alpha$.
It also does not justify identifying the constant term of a
characteristic power series with $p$: a finite $\Lambda$-module
has unit characteristic ideal, while it can have the nonzero
coinvariant group in (QI.2).  Passing from the actual ideal
$I_\alpha$ to a characteristic ideal discards that finite part.

The unevaluated arithmetic condition remains the unit norm index
in KT-ORDER-FIVE-GREENBERG-UNIT-GATE.  Under the surviving
$A_0^\alpha\simeq\mathbf F_p$ possibility, an index $p$ at some
layer would prove capitulation and bounded tower growth, but even
that is not itself triviality of the base class group.  No such
index or stronger annihilator is computed here.  The already
accepted periodic vanishing criterion and the actual base class
possibility are unchanged.  This is a control and scope check,
classified as NONMOVEMENT, not a proof that another Iwasawa
argument could never remove the remaining case.

\paragraph{Primary and carried source scope.}
The genus-field and principal-prime statements are the complete
proof of KT-ORDER-FIVE-GREENBERG-UNIT-GATE, reread for this
application; KT-ORDER-FIVE-STRICT-CYCLIC and
KT-ORDER-FIVE-PDIV-WEIGHT were also reread in full, as was the
existing abstract norm-model note.  The new proof above supplies
the compact inverse-limit and strict-character control steps.
The ordinary class-field inverse-limit convention is stated in
D.~C.~Marshall, \emph{Greenberg's conjecture and cyclotomic towers},
Acta Arithmetica \textbf{113.1} (2004), pp.~1--2.
Those pages and the complete statement of Theorem~1 on p.~3 were
read.  That theorem assumes Vandiver, a separate $\lambda_p=1$
condition, and a local-unit torsion bound for the full cyclotomic
field; it cannot be used to remove the present possible Vandiver
component.  No proof or conclusion from that theorem is imported.
Inflation--restriction, finite-group averaging, global class
field theory, and Pontryagin duality are foundational inputs.
No fresh proof of those general theories, Lean execution, or
new finite-degree computation is claimed.
